% ----------------------------------------------------------------
% AMS-LaTeX definitions for a Dissertation/Thesis ****************
% Template prepared by T.A. McWalter (March 2009)
% ----------------------------------------------------------------
\documentclass[a4paper,11pt,oneside]{report}

% PACKAGES -------------------------------------------------------

\usepackage{amsmath}
\usepackage{amsthm}
\usepackage{amsfonts}
\usepackage{amssymb}
\usepackage{fancyhdr}
\usepackage{footnote}
\usepackage{ifpdf}
\usepackage{mathrsfs}
\usepackage{multirow}
\usepackage[numbers]{natbib}
\usepackage{placeins}
\usepackage{xspace}

% GRAPHICS HANDLING ----------------------------------------------

\ifpdf
    \usepackage[pdftex]{graphicx}
    \DeclareGraphicsExtensions{.pdf, .jpg, .tif}
    \newcommand{\bb}{viewport}
\else
    \usepackage{graphicx}
    \DeclareGraphicsExtensions{.eps, .jpg}
    \newcommand{\bb}{bb}
\fi

% OTHER SET-UP ---------------------------------------------------

\pagestyle{fancy}
\lhead{\rightmark}
\chead{}
\rhead{\thepage}
\cfoot{}

\setlength{\topmargin}{0.5cm}
\setlength{\headsep}{1cm}
\setlength{\textwidth}{14.0cm}
\setlength{\headwidth}{14.0cm}
\setlength{\headheight}{14pt}
\setlength{\textheight}{22cm}
\setlength{\oddsidemargin}{1.3cm}
\setlength{\evensidemargin}{0cm}

\newcommand{\setlinespacing}[1]
           {\renewcommand{\baselinestretch}{#1}\small\normalsize}

% THEOREM STYLES -------------------------------------------------

\theoremstyle{plain}
\newtheorem{thm}{Theorem}[chapter]
\newtheorem{prop}[thm]{Proposition}
\newtheorem{cor}[thm]{Corollary}
\newtheorem{con}[thm]{Conjecture}
\newtheorem{lem}[thm]{Lemma}
\theoremstyle{definition}
\newtheorem{defn}[thm]{Definition}
\newtheorem{ass}[thm]{Assumption}
\newtheorem{exam}[thm]{Example}

% MACRO DEFINITIONS ----------------------------------------------
\def\n{\nonumber\\}
\def\h{\\\hline}
\newcommand{\theader}[1]{\multicolumn{1}{|c|}{\textbf{#1}}}

% Bond descriptions
\def\t{\ensuremath{t}\xspace}
\def\ti{\ensuremath{t_i}\xspace}
\def\T{\ensuremath{T}\xspace}
\def\N{\ensuremath{N}\xspace}
\def\R{\ensuremath{R}\xspace}
\def\C{\ensuremath{C_i}\xspace}
\def\Kp{\ensuremath{K^{p}_t}\xspace}
\def\Kc{\ensuremath{K^{c}_t}\xspace}
\def\Kv{\ensuremath{\kappa_t}\xspace}
\def\Tp{\ensuremath{\Omega^{p}}\xspace}
\def\Tc{\ensuremath{\Omega^{c}}\xspace}
\def\Tv{\ensuremath{\Omega^{v}}\xspace}
\def\ST{\ensuremath{S_T}\xspace}

% Statistical
\def\E[#1]{\Tilde{\mathbb{E}}\!\left[#1\right]}
\def\Et[#1]{\Tilde{\mathbb{E}}_t\!\left[#1\right]}
\def\Var[#1]{\Tilde{\mathbb{V}}\mathtt{ar}\!\left[#1\right]}
\def\Prob(#1){\Tilde{\mathbb{P}}\!\left(#1\right)}

% Derivatives
\def\S{\ensuremath{S_t}\xspace}
\def\Sdt{\ensuremath{S_{t + \dt}}\xspace}
\def\SDt{\ensuremath{S_{t + \Dt}}\xspace}
\def\W{\ensuremath{\tilde{W}_t}\xspace}
\def\q{\ensuremath{q_t}\xspace}
\def\qdt{\ensuremath{q_{t + \dt}}\xspace}
\def\dP{\ensuremath{d\Pi_t}\xspace}
\def\dS{\ensuremath{dS_t}\xspace}
\def\dt{\ensuremath{dt}\xspace}
\def\dW{\ensuremath{d\W}\xspace}
\def\dq{\ensuremath{d\q}\xspace}
\def\Vt{\frac{\partial V}{\partial t}}
\def\Vs{\frac{\partial V}{\partial S}}
\def\Vss{\frac{\partial^2 V}{\partial S^2}}

% Differences
\def\DS{\ensuremath{\delta S_t}\xspace}
\def\Dt{\ensuremath{\delta t}\xspace}
\def\DW{\ensuremath{\delta \W}\xspace}
\def\Dq{\ensuremath{\delta \q}\xspace}

% Finite difference
\def\Vjn{\ensuremath{V^{n}_j}\xspace}
\def\Vjnm{\ensuremath{V^{n - 1}_j}\xspace}
\def\Vjnp{\ensuremath{V^{n + 1}_j}\xspace}
\def\Vjmn{\ensuremath{V^{n}_{j - 1}}\xspace}
\def\Vjpn{\ensuremath{V^{n}_{j + 1}}\xspace}
\def\Vjmnm{\ensuremath{V^{n - 1}_{j - 1}}\xspace}
\def\Vjpnm{\ensuremath{V^{n - 1}_{j + 1}}\xspace}
\def\Vjmnp{\ensuremath{V^{n + 1}_{j - 1}}\xspace}
\def\Vjpnp{\ensuremath{V^{n + 1}_{j + 1}}\xspace}

\def\Von{\ensuremath{V^{n}_{0}\xspace}}
\def\Vomn{\ensuremath{V^{n}_{-1}\xspace}}
\def\Vopn{\ensuremath{V^{n}_{1}\xspace}}
\def\VJn{\ensuremath{V^{n}_{J}\xspace}}
\def\VJmn{\ensuremath{V^{n}_{J - 1}\xspace}}
\def\VJpn{\ensuremath{V^{n}_{J + 1}\xspace}}
\def\Vonm{\ensuremath{V^{n - 1}_{0}\xspace}}
\def\Vomnm{\ensuremath{V^{n - 1}_{-1}\xspace}}
\def\Vopnm{\ensuremath{V^{n - 1}_{1}\xspace}}
\def\VJnm{\ensuremath{V^{n - 1}_{J}\xspace}}
\def\VJmnm{\ensuremath{V^{n - 1}_{J - 1}\xspace}}
\def\VJpnm{\ensuremath{V^{n - 1}_{J + 1}\xspace}}

\def\Sz{\ensuremath{S_0}\xspace}
\def\Sj{\ensuremath{S_j}\xspace}
\def\SJ{\ensuremath{S_J}\xspace}
\def\Xj{\ensuremath{X_j}\xspace}

% Binomial model
\def\u{\ensuremath{u}\xspace}
\def\d{\ensuremath{d}\xspace}
\def\o{\ensuremath{o}\xspace}

\def\pu{\ensuremath{p_\u}\xspace}
\def\pd{\ensuremath{p_\d}\xspace}
\def\po{\ensuremath{p_\o}\xspace}

% Binomial model valuation
\def\V{\ensuremath{V_t}\xspace}
\def\Vu{\ensuremath{V^{u}_t}\xspace}
\def\Vd{\ensuremath{V^{d}_t}\xspace}
\def\X{\ensuremath{X_t}\xspace}
\def\P{\ensuremath{\Pi_t}\xspace}
\def\PDt{\ensuremath{\Pi_{t + \Dt}}\xspace}
\def\D{\ensuremath{\Delta_t}\xspace}

% Sundry
\def\th{\textsuperscript{th}\xspace}
\def\half{\tfrac{1}{2}}


% PREAMBLE -------------------------------------------------------

\title{Valuation of Convertible Bonds with Credit Risk}

\author{David Naylor\\[2cm]
Programme in Advanced Mathematics of Finance,\\
University of the Witwatersrand,\\
Johannesburg.\\[2cm]
\includegraphics[width=4cm]{../common/WitsColourLogo}\\[2cm]
}

% MAIN DOCUMENT --------------------------------------------------

\begin{document}

\maketitle \setlinespacing{1.20}

% DECLARATION ----------------------------------------------------

\def\abstractname{\Huge Declaration}
\begin{abstract}
I declare that this project is my own, unaided work. It is being submitted as partial fulfilment of the Degree of Bachelor of Science with Honours in the University of the Witwatersrand, Johannesburg. It has not been submitted before for any degree or examination in any other University.\\[3cm]
\noindent\rule{5cm}{0.5pt}\\[1cm]
\noindent \today
\end{abstract}

% ABSTRACT -------------------------------------------------------

\def\abstractname{\Huge Abstract}
\begin{abstract}
A convertible bond is a complex derivative that cannot be priced as a simple combination of bond and stock components.  Convertible bonds can be broken down as a bond with two embedded options (a put option for the investor and a call option for the issuer) and an option to convert the bond into stock.  Due to the multiple continuous options, the pricing of the convertible bond is path dependent.

This research project explores and implements a binary tree and finite difference scheme to price the convertible bond, taking into account credit risk.  Consideration is given to the sensitivities of the convertible bond and a possible hedging strategy.
\end{abstract}

% ACKNOWLEDGEMENTS -----------------------------------------------

\def\abstractname{\Huge Acknowledgements}
\begin{abstract}
I would like to thank my supervisor Tom McWalter for assigning me this very interesting problem and for the support he provided in understanding more difficult components of the project.

I would also like to thank Coenraad Labuschagne, the AMF Honours course coordinator and lecturer, and all the AMF Honours lecturers for providing me with the knowledge and understanding, both for an future and this project.

Lastly, I would like to thank my army of checkers and proof readers: Kay \mbox{McCallum}, Anne Naylor and Nigel Naylor.
\end{abstract}

% TABLE OF CONTENTS ----------------------------------------------

\setlinespacing{1}
\pagenumbering{roman}
\tableofcontents
\listoffigures
\listoftables
\clearpage
\setlinespacing{1.20}
\pagenumbering{arabic}
\setcounter{page}{1}

% =============================================================================

\chapter{Introduction}
This research project aims to price convertible bonds based on the paper by \citet{AFV03} and \citet{MK12}.

Convertible bonds are a hybrid instrument available on financial markets.  The convertible bond is an instrument that is similar to a normal bond, except the holder has the option to convert the bond into a specified number of shares.  The convertible bond typically also has embedded options whereby the issuer may buy back the convertible bond for a specified price and whereby the investor can force the issue to repurchase the convertible bond.  In the event of a default by the issuer the bond could have partial recovery or total default.

% -----------------------------------------------------------------------------

\section{Convertible Bonds with Credit Risk}
A convertible bond is an annuity, which pays coupons \C at \ti and pays \R at maturity \T, with the following embedded options:
\begin{itemize}
 \item \emph{A put option for the bond holder:} the bond holder may redeem the bond early for \Kp, provided $\t \in \Tp$.  This option is not applicable at termination, $T \not\in \Tp$.
 \item \emph{A call option for the bond issuer:} the bond issuer may force early redemption of the bond for \Kc, provided $\t \in \Tc$.  Normally $\Kc > \Kp$.  This option is not applicable at termination, $T \not\in \Tc$.
 \item \emph{A conversion option:} the bond holder may convert the bond into \Kv shares, for a value of $\Kv\S$, provided $\t \in \Tv$.  To optimise profits the bond holder would rather convert than exercise if $\Kp < \Kv\S$.  This option supersedes the call option (i.e. if the bond issuer forces early redemption the bond holder may still convert the bond into shared).  This option is applicable at termination, $T \in \Tv$.
\end{itemize}

In the event of default the bond holder may either recover the remains of the annuity or convert into the stock, which ever is worth more.  Default is considered terminal.

Suppose the recovery of the bond is $\gamma$, and the drop in share price is $\eta$, then the value of the convertible bond in default is\footnote{This is assuming the bond recovery is based on the redemption value.  Another possibility is to base the recovery on the bond value at time \t}:
\begin{equation}
 \X = \max(\gamma\R, (1 - \eta)\Kv\S)
\end{equation}


Table \ref{tab:choices} illustrated the various choices made by the bond holder and the bond issuer, and under what circumstances those choices are made.  Figure \ref{fig:choices} illustrates the actions taken on the bond over the lifetime of a convertible bond.  The impact of coupons is left as an implementation detail.
\begin{table}[ht]
 \centering
 \begin{tabular}{|l|l|l|}\hline
  \theader{Action}    & \theader{Payoff} & \theader{Condition} \h
  Put                 & \Kp              & $(\V \le \Kp)\wedge(\t \in \Tp)\wedge[(\Kv\S < \V)\lor(\t \not\in \Tv)]$ \h
  Call                & \Kc              & $(\V \ge \Kc)\wedge(\t \in \Tc)\wedge[(\Kv\S < \Kc)\lor(\t \not\in\Tv)]$ \h
  Conversion          & $\Kv\S$          & $(\Kv\S \ge \V)\wedge(\t \in \Tv)$ \h
  Forced conversion   & $\Kv\S$          & $(\V > \Kv\S \ge \Kc)\wedge(\t \in \Tv)\wedge(\t \in \Tc)$ \h
  Redemption          & \R               & $(\t = \T)\wedge[(\Kv\S \le \R)\lor(\t \not\in \Tv)]$ \h
  Hold                &                  & \emph{otherwise} \h
 \end{tabular}
 \caption[Payoff for the convertible bond]{Payoff for the convertible bond.  \V is assumed to be the intrinsic value in this context only}
 \label{tab:choices}
\end{table}

% -----------------------------------------------------------------------------

\section{Literature Review}
\citet{I77} started the literature of pricing a convertible bond with extensions from \citet{BS77} and \citet{BS80}.  The original approach was to treat the bond and equity as components of the issuer's value and to treat default as when the issuer's value drops below a point where it can no longer meet its financial obligations.  An overview of this type of approach is provided by \citet{N96} and criticisms are addressed by \citet{JT95}.  The main problems with this model is that the issuer's value is not directly observable, difficult to parameterise and all senior debt to the convertible bond also needs to be priced.

A second approach was to price the convertible bond based on the issuer's stock price \cite{MS86}.  A refined method, called ``reduced form'', treats default as a discrete jump in time.  \Citet{AFV03} was the first to allow stock jumps that didn't go to zero.  The probability of the loss jump over a short period of time is described by a hazard rate.

Also some Monte Carlo based pricing methods have been considered as proposed by \citet{B89} with \citet{G03} providing an optimisation approach to handle optimal early exercise of the American options.  Further improvements to the pricing methods were done by \citet{AKW08}.

Other methods used for pricing include a finite element method by \citet{BBH03} and a binomial tree method by \citet{TKN01} and \citet{AFV02}.

% =============================================================================

\chapter{Finite Difference Model}
Describe the finite difference model based on the paper by \citet{AFV03}.

% -----------------------------------------------------------------------------

\section{Stochastic Process}
Consider a stock with the following stochastic model for stock price movements:
\begin{equation}
 \dS = \mu\S\dt + \sigma\S\dW - \eta\S\dq
\end{equation}
If we impose that the risk neutral expectation that the stock price evolution is the risk free rate, $\E[\dS] = r\S\dt$, then $\mu = (r + \lambda\eta)$ and the differential equation becomes:
\begin{equation}
 \dS = (r + \lambda\eta)\S\dt + \sigma\S\dW - \eta\S\dq
\end{equation}
where:

\begin{tabular}{ll}
 $\mu$          & is the drift rate of the stock price \\
 $r$            & is the risk free rate \\
 $\sigma$       & is the log-volatility of the stock price \\
 $\lambda$      & is the hazard rate \\
 $\eta$         & is the percentage drop in stock price on a default event \\
 \dW            & is a Wiener process \\
 \dq            & is a Poisson jump process where the first jump is the default event \\
\end{tabular}

% -----------------------------------------------------------------------------

\section{Derivation}
Consider a portfolio, \P, that has a long position in the derivative, \V, and a short position of \D in the stock, \S.  The residual $\P = \V - \D\S$ is invested at the risk neutral rate.  This yields a differential equation for \P:
\begin{align}
 \dP &= r\P\dt \n
     &= r(\V - \D\S)\dt \label{eq:rPdt}
\end{align}

In the event of default the portfolio holdings $\P^x$ has a long position in the derivative with residual value \X and a short position of \D in the stock with residual value of $(1 - \eta)\S$.  Thus the difference of the portfolio holdings when going into default is:
\begin{align}
 \dP^x &= \P^x - \P \n
       &= \X - (\V - \eta\D\S) \label{eq:dPx}
\end{align}

At any given instant the portfolio can either default or evolve in the normal stochastic manner.  The probability of default in any given instant is $\lambda\dt$ and the probability of non-default is $(1 - \lambda)\dt$.  Taking the expectation of the differential of \P and applying Ito's lemma to the stochastic component:
\begin{align}
 \Et[\dP] &= (1 - \lambda\dt)[\left(\Vt + \mu\S\Vs + \half\sigma^2\S^2\Vss - \mu\S\D\right)\dt + \n
          &\qquad \left(\sigma\S\Vs - \sigma\S\D\right)\dW] + \lambda\dt(\X - (\V - \eta\D\S))
\end{align}
By choosing \D to eliminate the \dW term, $\D = \Vs$, the differential simplifies to:
\begin{equation}
 \Et[\dP] = \left(\Vt + \half\sigma^2\S^2\Vss\right)\dt + \lambda\dt(\X - (\V - \eta\D\S)) \label{eq:EdP}
\end{equation}

Equating the differential \eqref{eq:rPdt} and the expected differential \eqref{eq:EdP} for \P the partial differential equation for a derivative with credit risk is derived:
\begin{align}
 r(\V - \D\S)\dt &= \left(\Vt + \half\sigma^2\S^2\Vss\right)\dt + \lambda\dt(\X - (\V - \eta\D\S)) \n
 (r + \lambda)\V &= \Vt + (r + \lambda\eta)\S\Vs + \half\sigma^2\S^2\Vss + \lambda\X
\end{align}

% -----------------------------------------------------------------------------

\section{Finite Difference}
The derived PDE can be solved numerically by constructing a lattice of stock prices and time.  \Vjn will denote the lattice value at the j\th stock node and n\th time node where the step between stock nodes is \DS and the step between time nodes is \Dt.  \Sj will denote the stock value at the j\th stock node and \Xj will denote the defaulted value at the j\th stock node.  The stock prices go from a minimum of \Sz to a maximum of \SJ.

All finite difference equations derived here can be expressed as functions of $\alpha$, $\beta$ and $\kappa$ where:
\begin{align*}
 \alpha &= \half\sigma^2\Sj^2\frac{\Dt}{\DS^2} - (r + \lambda\eta)\Sj\frac{\Dt}{2\DS} \\
 \beta  &= - (r + \lambda) - \sigma^2\Sj^2\frac{\Dt}{\DS^2} \\
 \kappa &= \half\sigma^2\Sj^2\frac{\Dt}{\DS^2} + (r + \lambda\eta)\Sj\frac{\Dt}{2\DS} \\
\end{align*}
and it is assumed that for boundary nodes:
\begin{align*}
 \frac{\Von - \Vomn}{\DS} &= \frac{\Vopn - \Von}{\DS} \\
                    \Vomn &= 2\Von - \Vopn \\
 \frac{\VJn - \VJmn}{\DS} &= \frac{\VJpn - \VJn}{\DS} \\
                    \VJpn &= 2\VJn - \VJmn \\
\end{align*}

% -----------------------------------------------------------------------------

\subsection{Explicit Scheme}
Using the backwards difference equation for the temporal derivative the finite difference equation is:
\begin{displaymath}
 (r + \lambda)\Vjn = \frac{\Vjn - \Vjnm}{\Dt} + (r + \lambda\eta)\Sj\frac{\Vjpn - \Vjmn}{2\DS} + \half\sigma^2\Sj^2\frac{\Vjmn - 2\Vjn + \Vjpn}{\DS^2} + \lambda\Xj
\end{displaymath}
\begin{align}
 \Vjnm &= \alpha\Vjmn + (1 + \beta)\Vjn + \kappa\Vjpn + \lambda\Dt\Xj \\
 \Vonm &= (1 + \beta + 2\alpha)\Von + (\kappa - \alpha)\Vopn + \lambda\Dt\Xj \\
 \VJnm &= (\alpha - \kappa)\Vjmn + (1 + \beta - 2\kappa)\VJn + \lambda\Dt\Xj
\end{align}

% -----------------------------------------------------------------------------

\subsection{Implicit Scheme}
Using the backwards difference equation for the temporal derivative and evaluating the spacial differentials at the previous time step the finite difference equation is:
\small\begin{displaymath}
 (r + \lambda)\Vjnm = \frac{\Vjn - \Vjnm}{\Dt} + (r + \lambda\eta)\Sj\frac{\Vjpnm - \Vjmnm}{2\DS} + \half\sigma^2\Sj^2\frac{\Vjmnm - 2\Vjnm + \Vjpnm}{\DS^2} + \lambda\Xj
\end{displaymath}\normalsize
\begin{align}
 -\alpha\Vjmnm + (1 - \beta)\Vjnm - \kappa\Vjpnm &= \Vjn + \lambda\Dt\X \\
 (1 - \beta - 2\alpha)\Vonm - (\kappa - \alpha)\Vopnm &= \Von + \lambda\Dt\X \\
 -(\alpha - \kappa)\VJmnm + (1 - \beta - 2\kappa)\VJnm &= \VJn + \lambda\Dt\X
\end{align}

% -----------------------------------------------------------------------------

\subsection{Crank-Nicolson Scheme}
Using the backwards difference equation for the temporal derivative and taking the average of the spacial differentials at the current and previous time step the finite difference equation is:
\small\begin{align*}
 (r + \lambda)\half(\Vjn + \Vjnm) &= \frac{\Vjn - \Vjnm}{\Dt} + (r + \lambda\eta)\Sj\frac{1}{2}\left(\frac{\Vjpn - \Vjmn}{2\DS} + \frac{\Vjpnm - \Vjmnm}{2\DS}\right) + \\
  &\quad \half\sigma^2\Sj^2\frac{1}{2}\left(\frac{\Vjmn - 2\Vjn + \Vjpn}{\DS^2} + \frac{\Vjmnm - 2\Vjnm + \Vjpnm}{\DS^2}\right) + \lambda\Xj
\end{align*}
\begin{align}
 -\alpha\Vjmnm + (2 - \beta)\Vjnm - \kappa\Vjpnm &= \alpha\Vjmn + (2 + \beta)\Vjn + \kappa\Vjpn + 2\lambda\Dt\X \\
 (2 - \beta - 2\alpha)\Vonm - (\kappa - \alpha)\Vopnm &= (2 + \beta + 2\alpha)\Von + (\kappa + \alpha)\Vopn + 2\lambda\Dt\X \\
 -(\alpha - \kappa)\VJmnm + (2 - \beta - 2\kappa)\VJnm &= (\alpha - \kappa)\VJmn + (2 + \beta + 2\kappa)\VJn + 2\lambda\Dt\X
\end{align}\normalsize

% -----------------------------------------------------------------------------

\section{Coupons}
TODO: figure out how to handle coupons

% =============================================================================

\chapter{Binomial Model}
The Binomial Model with Credit Risk was first derived by \citet{MK12}, who showed that this model converges, in continuous time, to the \citet{AFV03} model.  \Citet{MK12} also showed the valuation method of this model is the same as in the classical binomial model.

TODO: add further explanation and references.

% -----------------------------------------------------------------------------

\section{Binomial Process}
Consider a binomial model that has time step \Dt, up and down steps, \u and \d with probability \pu and \pd\footnote{$\pd = 1 - \pu - \po = e^{-\lambda\Dt} - \pu$} respectively, and where the probability of default is \po\footnote{The time till the first jump follows an exponential distribution with intensity $\lambda$ and has probability $\Prob(q_{t + \dt} = 1 | q_t = 0) = \po = 1 - e^{-\lambda\Dt}$}.  The mean and variance of this model is: (TODO: insert picture illustration)
\begin{align}
 \E[\frac{\SDt}{\S}]   &= \u\pu + \d\pd + (1 - \eta)\po \n
                       &= \u\pu + \d(e^{-\lambda\Dt} - \pu) + (1 - \eta)(1 - e^{-\lambda\Dt}) \label{eq:BM_E} \\
 \Var[\frac{\SDt}{\S}] &= \u^2\pu + \d^2\pd + (1 - \eta)^2\po - \E[\frac{\SDt}{\S}]^2 \n
                       &= \u^2\pu + \d^2(e^{-\lambda\Dt} - \pu) + (1 - \eta)^2(1 - e^{-\lambda\Dt}) - \E[\frac{\SDt}{\S}]^2 \label{eq:BM_Var}
\end{align}

The mean and variance of the stock's stochastic differential equation is:
\begin{align}
 \E[\frac{\dS}{\S}]    &= r\dt \\
 \E[\frac{\Sdt}{\S}]   &= \E[\frac{\S}{\S} + \frac{\dS}{\S}] \n
                       &= 1 + \E[\frac{\dS}{\S}] \label{eq:SM_E} \\
 \Var[\frac{\dS}{\S}]  &= (\sigma^2 + \lambda\eta^2)\dt \\
 \Var[\frac{\Sdt}{\S}] &= \Var[\frac{\S}{\S} + \frac{\dS}{\S}] \n
                       &= \Var[\frac{\dS}{\S}] \label{eq:SM_Var}
\end{align}


% -----------------------------------------------------------------------------

\section{Derivation}
If one equates the first moment of the stochastic model \eqref{eq:SM_E}\footnote{$\E[\frac{\SDt}{\S}] = e^{r\Dt}$}, in \Dt time, with that of the binomial model \eqref{eq:BM_E} then:
\begin{align}
 e^{r\Dt}       &= \u\pu + \d(e^{-\lambda\Dt} - \pu) + (1 - \eta)(1 - e^{-\lambda\Dt}) \n
 \pu(\u - \d)   &= e^{r\Dt} - \d e^{-\lambda\Dt} - (1 - \eta)(1 - e^{-\lambda\Dt}) \label{eq:p(u-d)}\\
 \pu            &= \frac{e^{r\Dt} - \d e^{-\lambda\Dt} - (1 - \eta)(1 - e^{-\lambda\Dt})}{\u - \d} \\
 \therefore\quad\pd &= -\frac{e^{r\Dt} - \u e^{-\lambda\Dt} - (1 - \eta)(1 - e^{-\lambda\Dt})}{\u - \d}
\end{align}
and if one equate the second moment about the mean of the stochastic model \eqref{eq:SM_Var}, in \Dt time, with that of the binomial model \eqref{eq:BM_Var} then\footnote{Substituting equation \eqref{eq:p(u-d)} into \eqref{eq:Var=1} one gets \eqref{eq:Var=2}}:
\begin{align}
 (\sigma^2 + \lambda\eta^2)\Dt &= \u^2\pu + \d^2(e^{-\lambda\Dt} - \pu) + (1 - \eta)^2(1 - e^{-\lambda\Dt}) - \E[\frac{\SDt}{\S}]^2 \n
                        &= (\u^2 - \d^2)\pu + \d^2 e^{-\lambda\Dt} +  (1 - \eta)^2(1 - e^{-\lambda\Dt}) - e^{-2r\Dt} \label{eq:Var=1} \\
                        &= (\u + \d)(e^{r\Dt} - \d e^{-\lambda\Dt} - (1 - \eta)(1 - e^{-\lambda\Dt})) \label{eq:Var=2} \\
                        &\qquad + \d^2 e^{-\lambda\Dt} + (1 - \eta)^2(1 - e^{-\lambda\Dt}) - e^{-2r\Dt} \n
                        &= (\u + \d)(e^{r\Dt} - (1 - \eta)(1 - e^{-\lambda\Dt})) - \u\d e^{-\lambda\Dt} \label{eq:Var=} \\
                        &\qquad + (1 - \eta)^2(1 - e^{-\lambda\Dt}) - e^{-2r\Dt} \nonumber
\end{align}
If one assumes $\Dt^2 = 0$, $\u\d = 1$ and $\u = e^{\sqrt{A\Dt}}$ and the Taylor series expansion is taken for all exponential terms, then:
\begin{align}
 \u                     &= 1 + \sqrt{A\Dt} + \frac{A\Dt}{2!} + \frac{(A\Dt)^{\frac{3}{2}}}{3!} \n
 \d                     &= 1 - \sqrt{A\Dt} + \frac{A\Dt}{2!} - \frac{(A\Dt)^{\frac{3}{2}}}{3!} \n
 \u + \d                &= 2 + A\Dt \label{eq:u+d}
\end{align}
and substituting \eqref{eq:u+d} into \eqref{eq:Var=} with the appropriate expansions then:
\begin{align}
 (\sigma^2 + \lambda\eta^2)\Dt &= (2 + A\Dt)(1 + r\Dt - \lambda\Dt(1 - \eta)) - (1 - \lambda\Dt) \n
                        &\qquad + \lambda\Dt(1 - \eta)^2 - (1 - 2r\Dt) \n
                        &= A\Dt + \lambda\Dt - \lambda\Dt(1 - \eta)(1 + \eta) \n
                        &= A\Dt + \lambda\eta^2\Dt \n
 A\Dt                   &= \sigma^2\Dt
\end{align}

% -----------------------------------------------------------------------------

\section{Restrictions}
The time step \Dt is assumed to be strictly positive ($\Dt > 0$), as time is monotonically increasing.

Although $\sigma$ could be positive or negative, the formul\ae\ use $\sigma$ where, strictly, $|\sigma|$ should be used.  If $\sigma = 0$ then the model is no longer useful.  The restriction $\sigma > 0$ may be introduced for convenience and without loss of generality.

An implicit assumption this model makes is that \pu, \pd and \po are a valid probability measure.  This imposes the following conditions:
\begin{align}
 \min(\pu, \pd, \po) &\ge 0 \\
 \pu + \pd + \po &= 1 \label{eq:bin_prob_limit}
\end{align}

Consider \po, the probability of default.  Its value is equal to the cumulative density function of an exponential distribution with intensity $\lambda$.  As such \po is already confined to the interval $[0, 1]$ with the restriction that $\lambda > 0$.  If $\lambda = 0$, no default is possible, then \po will have a value of 0 and is a valid probability thus the restriction on $\lambda$ can be relaxed to $\lambda \ge 0$.

Consider \pd, the probability of a down movement.  Its value is defined in terms of \pu and \po.  If \pu and \po are valid probabilities, within bounds, then so will \pd.

Consider \pu, the probability of an up movement.  Taking into account the definition of \pd, this requires that \pu satisfies:
\begin{align}
 0 &\le \pu \le e^{-\lambda \Dt} \\
 0 &\le \frac{e^{r\Dt} - \d e^{-\lambda\Dt} - (1 - \eta)(1 - e^{-\lambda\Dt})}{\u - \d} \le e^{-\lambda \Dt} \n
 0 &\le e^{r\Dt} - \d e^{-\lambda\Dt} - (1 - \eta)(1 - e^{-\lambda\Dt}) \le e^{-\lambda \Dt}(\u - \d) \n
 e^{-\lambda\Dt}(\d - (1 - \eta)) &\le e^{r\Dt} - (1 - \eta) \le e^{-\lambda\Dt}(\u - (1 - \eta)) \n
 \ln\left(\frac{d - (1 - \eta)}{e^{r\Dt} - (1 - \eta)}\right) &\le \lambda\Dt \le \ln\left(\frac{u - (1 - \eta)}{e^{r\Dt} - (1 - \eta)}\right) \\
 \Dt &\le \frac{1}{\lambda}\ln\left(\frac{u - (1 - \eta)}{e^{r\Dt} - (1 - \eta)}\right)
\end{align}
Using the limit on both \Dt and $\lambda$ the inequalities around the $\lambda\Dt$ are also limited thus:
\begin{align}
 0 &\le \ln\left(\frac{u - (1 - \eta)}{e^{r\Dt} - (1 - \eta)}\right) \n
 e^{r\Dt} - (1 - \eta) &\le u - (1 - \eta) \n
 r\Dt &\le \sigma \sqrt{\Dt} \n
 \Dt &\le \frac{\sigma^2}{r^2}
\end{align}
The inequality of $\ln\left(\frac{d - (1 - \eta)}{e^{r\Dt} - (1 - \eta)}\right) \le 0$ results in the same inequality as above.

Based on the probability limits \eqref{eq:bin_prob_limit} no restriction is placed on $r$, the risk free rate, or $\eta$, the percentage drop in stock price.

% -----------------------------------------------------------------------------

\section{Valuation}
TODO: add some fluff

\subsection{Valuation without Coupons}
Consider a portfolio, \P, that has a long position in the derivative, \V, and a short position of \D in the stock, \S.  The residual $\P = \V - \D\S$ is invested at the risk neutral rate.  At time $t + \Dt$ one should have:
\begin{equation*}
 \PDt =
 \left\{\begin{array}{l@{\qquad\text{with probability }}l}
  \Vu - \D\S u          & \pu \\
  \Vd - \D\S d          & \pd \\
  \X - (1 - \eta)\D\S   & \po \\
 \end{array}\right.
\end{equation*}
where

\begin{tabular}{ll}
 $\Vu$          & is the value of the portfolio on an up-step at time $t + \Dt$ \\
 $\Vd$          & is the value of the portfolio on a down-step at time $t + \Dt$ \\
 $\X$           & is the value of default at time $t$ \\
\end{tabular}

If one wishes to hedge against up and down movements of the stock then:
\begin{align}
    \Vu - \D\S u &= \Vd - \D\S d \n
     \D\S(d - u) &= \Vd - \Vu \n
              \D &= \frac{\Vu - \Vd}{\S(u - d)}
\end{align}

The required value of \P at time $t + \Dt$ is:
\begin{align}
 \PDt  &= \P e^{r\Dt} \n
       &= \left(\V - \frac{\Vu - \Vd}{u - d}\right) e^{r\Dt} \label{eq:P}
\end{align}
as \P is invested at the risk free rate, and taking the expected value of $\PDt$, which must equal \eqref{eq:P}, one gets:
\begin{align}
 \left(\V - \frac{\Vu - \Vd}{u - d}\right) e^{r\Dt} &= \Et[\PDt] \n
        &= \left(\Vu - \frac{\Vu - \Vd}{u - d}u\right) e^{-\lambda\Dt} \n
        &\qquad + \left(\X - \frac{\Vu - \Vd}{u - d}(1 - \eta)\right)(1 - e^{-\lambda\Dt}) \n
        &= \Vu\left(e^{-\lambda\Dt} - \frac{ue^{-\lambda\Dt} + (1 - \eta)(1 - e^{-\lambda\Dt})}{u - d}\right) \n
        &\qquad + \Vd\left(\frac{ue^{-\lambda\Dt} + (1 - \eta)(1 - e^{-\lambda\Dt})}{u - d}\right) + \X(1 - e^{-\lambda\Dt}) \n
 \V e^{r\Dt} &= \Vu\frac{e^{r\Dt} - de^{-\lambda\Dt} - (1 - \eta)(1 - e^{-\lambda\Dt})}{u - d} \n
        &\qquad + \Vd\frac{ue^{-\lambda\Dt} + (1 - \eta)(1 - e^{-\lambda\Dt}) - e^{r\Dt}}{u - d} + \X(1 - e^{-\lambda\Dt}) \n
 \V     &= e^{-r\Dt}(\Vu\pu + \Vd\pd + \X\po)
\end{align}
thus the pricing method of this binomial model is the same as the pricing method of the classical binomial model:
\begin{align}
 V_t    &= e^{-r\Dt}(\Vu\pu + \Vd\pd + \X\po) \n
        &= e^{-r\Dt}\Et[V_{t + \Dt}]
\end{align}

% -----------------------------------------------------------------------------

\subsection{Valuation with Coupons}
Consider a portfolio, as above, with the add of coupon cash flows.  At time $t + \Dt$ one should have:
\begin{equation*}
 \PDt =
 \begin{cases}
  \Vu - \D\S u + c_i e^{t + \Dt - t^{c}_i}      & \text{with probability } \pu \\
  \Vd - \D\S d + c_i e^{t + \Dt - t^{c}_i}      & \text{with probability } \pd\\
  \X - (1 - \eta)\S + c_i (1 - q_{t^{c}_i})e^{t + \Dt - t^{c}_i} & \text{with probability } \po\\
 \end{cases}
\end{equation*}
where

\begin{savenotes}
\begin{tabular}{ll}
 $c_i$          & is the i\textsuperscript{th} coupon value \\
 $t^{c}_i$      & is the time of the i\textsuperscript{th} coupon, with $t^{c}_i \in (t, t + \Dt)$ \\
 $q_{t^{c}_i}$  & is an indicator variable that default happens after coupon payment\footnote{$\Prob(q_{t^{c}_i} = 0 | q_s = 1, s \in (t, t + \Dt]) = \frac{e^{-\lambda(t^{c}_i - t)} - e^{-\lambda\Dt}}{1 - e^{-\lambda\Dt}}$, the probability that the coupon arrives before default occurs, given that default occurs within the interval}
\end{tabular}
\end{savenotes}

Using the above method it is trivial to show that:
\begin{align}
 V_t    &= e^{-r\Dt}(\Vu\pu + \Vd\pd + \X\po) + c_i e^{-(r + \lambda)(t^{c}_i - t)} \n
        &= e^{-r\Dt}\Et[V_{t + \Dt}] + c_i e^{-(r + \lambda)(t^{c}_i - t)}
\end{align}
if, however, $t^{c}_i = t$, thus $c_i$ arrives with certainty at $t$, it is trivially shown that:
\begin{align}
 V_t    &= e^{-r\Dt}(\Vu\pu + \Vd\pd + \X\po) + c_i \n
        &= e^{-r\Dt}\Et[V_{t + \Dt}] + c_i
\end{align}

% -----------------------------------------------------------------------------

\section{Synthesising and Hedging Strategy}
Consider the strategy where $\V - \D\S$ is invested at the risk free rate and $\D\S$ is held in stock.  This creates the portfolio $\P = \D\S + (\V - \D\S) = \V$ and in one time step the portfolio is worth:
\begin{equation}
 \PDt = \left\{\begin{array}{l@{\qquad\text{with probability }}l}
                \D\S\u + e^{r\Dt}(\V - \D\S)         & \pu \\
                \D\S\d + e^{r\Dt}(\V - \D\S)         & \pd \\
                (1 - \eta)\D\S + e^{r\Dt}(\V - \D\S) & \po
               \end{array}\right. \label{eq:PDt}
\end{equation}

Now consider the portfolio value for an up-movement in stock:
\begin{align}
 \PDt^u &= \D\S\u + e^{r\Dt}(\V - \D\S) \n
        &= \D\S(\u - e^{r\Dt}) + e^{r\Dt}\V \n
        &= (\Vu - \Vd)\left(\frac{\u - e^{r\Dt}}{u - d}\right) + (\Vu\pu + \Vd\pd + \X\po) \n
%
%        &= (\Vu - \Vd)\frac{\u - e^{r\Dt}}{u - d} + \left(\Vu\left(\frac{e^{r\Dt} - \d e^{-\lambda\Dt} - (1 - \eta)\po}{\u - \d}\right)\right. \n
        &\qquad \left.+ \Vd\left(\frac{\u e^{-\lambda\Dt} + (1 - \eta)\po - e^{r\Dt}}{\u - \d}\right) + \X\po\right) \n
%
        &= \Vu + \left(\Vu\left(\frac{\d - \d e^{-\lambda\Dt} - (1 - \eta)\po}{\u - \d}\right)\right. \n
        &\qquad \left.+ \Vd\left(\frac{\u e^{-\lambda\Dt} + (1 - \eta)\po - \u}{\u - \d}\right) + \X\po\right) \n
%
        &= \Vu + \left(\left(\frac{\Vu\d - \Vd\u}{\u - \d}\right) - (1 - \eta)\left(\frac{\Vu - \Vd}{\u - \d}\right) + \X\right)\po \n
        &= \Vu + \left(-\Vu + \left(\frac{\Vu - \Vd}{\u - \d}\right)\u - (1 - \eta)\D\S + \X\right)\po \n
        &= \Vu + (\X - (\Vu  - \D\S\u) - (1 - \eta)\D\S)\po
\end{align}

Similarly the portfolio value for a down-movement in stock:
\begin{equation}
 \PDt^d = \Vd + (\X - (\Vd  - \D\S\d) - (1 - \eta)\D\S)\po
\end{equation}

The portfolio value for a default in stock:
\begin{align}
 \PDt ^o &= (1 - \eta)\D\S + e^{r\Dt}(\V - \D\S) \n
         &= \D\S((1 - \eta) - e^{r\Dt}) + e^{r\Dt}\V \n
         &= (\Vu - \Vd)\left(\frac{(1 - \eta) - e^{r\Dt}}{u - d}\right) + (\Vu\pu + \Vd\pd + \X\po) \n
         &= \X\po - \Vu\left(\frac{\d e^{-\lambda\Dt} - (1 - \eta)e^{-\lambda\Dt}}{u - d}\right) + \Vd\left(\frac{\u e^{-\lambda\Dt} - (1 - \eta)e^{-\lambda\Dt}}{u - d}\right) \n
         &= \X\po - \left(\frac{\Vu\d - \Vd\u}{u - d}\right)e^{-\lambda\Dt} + (1 - \eta)\left(\frac{\Vu - \Vd}{u - d}\right)e^{-\lambda\Dt} \n
         &= \X\po + ((\Vu - \D\S\u) + (1 - \eta)\D\S)e^{-\lambda\Dt} \n
         &= \X - (\X - (\Vu - \D\S\u) - (1 - \eta)\D\S)e^{-\lambda\Dt} \n
\end{align}

Consider $\H = \X - (\Vu - \D\S\u) - (1 - \eta)\D\S$, which could be considered the un-adjusted\footnote{Adjustments need to be made for the probability of default} default cost after hedging.  $\H\po$ is received for both an up- and down- step in stock price, so $\H\po$ is received $e^{-\lambda\Dt}$ of the time and $\H e^{-\lambda\Dt}$ is paid \po of the time, on balance.  Each time step effectively insures default for that time step by receiving on non-default events what needs to be paid on a default event.

An important note about using the self-financing property of this model in practise is that it will require a large pool of uncorrelated portfolios with the same adjusted default cost after hedging and actuarial consideration to effectively handle the ``premiums'' and reserves.

Equation \ref{eq:PDt} can be rewritten as:
\begin{equation}
 \PDt = \left\{\begin{array}{l@{\qquad\text{with probability }}l}
                \Vu + \H\po             & \pu \\
                \Vd + \H\po             & \pd \\
                \X - \H e^{-\lambda\Dt} & \po
               \end{array}\right.\label{eq:PDt2}
\end{equation}

It is trivial to confirm that the expected value of \PDt is $e^{r\Dt}\P$.

% -----------------------------------------------------------------------------

\section{Parameters and Formul\ae\ of Model}
Based on the above Table \ref{tab:bin_parameter} lists the parameters required for, Table \ref{tab:bin_limit} lists the limits of, and Table \ref{tab:bin_formulae} lists the formul\ae\ to price, this binomial tree.

\begin{table}[ht]
 \centering
 \begin{tabular}{|r|l|}\hline
  \theader{Parameter} & \theader{Description} \h
  $r$                 & Risk free rate \h
  $\sigma$            & Log-volatility of the stock price \h
  $\lambda$           & Hazard rate of default \h
  $\eta$              & Percentage drop of stock price in a default event \h
  \Dt                 & Time-step \h
  \Vu                 & Value of portfolio for an up-step at $t + \Dt$ \h
  \Vd                 & Value of portfolio for a down-step at $t + \Dt$ \h
  \X                  & Value of defaulted portfolio at $t$ \h
  $c_i$               & Value of the i\textsuperscript{th} coupon payment \h
  $t^{c}_i$           & Time of the i\textsuperscript{th} coupon payment \h
 \end{tabular}
 \caption{Parameters of the Binomial Model with Credit Risk}
 \label{tab:bin_parameter}
\end{table}

\begin{savenotes}
\begin{table}[ht]
 \centering
 \begin{tabular}{|c|l|}\hline
  \theader{Limit}                & \theader{Description} \h
  $0 < \sigma$                   & Volatility must be positive \h
  $0 \le \lambda$                & Hazard rate must be non-negative \h
  $0 < \Dt$                      & Time step must be positive \h
  $\Dt \le \frac{\sigma^2}{r^2}$ & Time step must be small enough to handle volatility\footnote{If $r = 0$ then this inequality is not applicable.} \h
  $\Dt \le \frac{1}{\lambda}\ln\left(\frac{u - (1 - \eta)}{e^{r\Dt} - (1 - \eta)}\right)$ & Time step must be small enough to handle hazard rate\footnote{If $\lambda = 0$ then this inequality is not applicable.} \h
 \end{tabular}
 \caption{Limits of the Binomial Model with Credit Risk}
 \label{tab:bin_limit}
\end{table}
\end{savenotes}

\begin{table}[ht]
 \centering
 \begin{tabular}{|r@{ = }l|l|l|}\hline
  \multicolumn{2}{|c|}{\textbf{Formul\ae}} & \theader{Description} \h
  \u     & $e^{\sigma\sqrt{\Dt}}$ &
     Multiplier for an up-step \h
  \d     & $e^{-\sigma\sqrt{\Dt}}$ &
     Multiplier for a down-step \h
  \pu    & $\frac{e^{r\Dt} - \d e^{-\lambda\Dt} - (1 - \eta)(1 - e^{-\lambda\Dt})}{\u - \d}$ &
     Probability of an up-step \h
  \pd    & $e^{-\lambda\Dt} - \pu$ &
     Probability of a down-step \h
  \po    & $1 - e^{-\lambda\Dt}$ &
     Probability of default \h
  \V     & $e^{-r\Dt}(\Vu\pu + \Vd\pd + \X\po) + c_i$ &
     Value of portfolio at $t$ for $t^{c}_i = t$ \h
  \D     & $\tfrac{\Vu - \Vd}{S(\u - \d)}$ &
     Delta ratio for hedging \h
  \H     & $\X - (\Vu - \D\S\u) - (1 - \eta)\D\S$ &
     Unadjusted default cost after hedging \h
 \end{tabular}
 \caption{Formal\ae\ for the Binomial Model with Credit Risk}
 \label{tab:bin_formulae}
\end{table}

% =============================================================================

\chapter{Numerical Example}
In the following analysis of the numerical results Table \ref{tab:bond_parameters}, the bond parameters, and Table \ref{tab:stock_parameters}, the stock parameters, are used as a basis.  There are three different types of stock price decreases that are used: total default when the stock losses all value, typical default when the stock losses 40\% and partial default where the stock price doesn't lose any value.

If $\t^+$ is the time immediately after a coupon payment and $\t^-$ is the immediately before a coupon payment then the coupon payment is handled as follows:
\begin{equation}
 V_{t^{-}_i} = V_{t^{+}_i} + \C
\end{equation}
and it is assumed that embedded options can only be exercised after receiving the coupon payment (i.e. at $t^{+}_i$).

\begin{savenotes}
\begin{table}[ht]
 \centering
 \begin{tabular}{|c|l|c|} \hline
  \theader{Component}         & \theader{Parameter} & \theader{Value} \h
  \multirow{5}{*}{Annuity}    & Notional\footnote{$\R := 104$} & $100$ \\
                              & Coupon\footnote{$\C := 4$} & 8\% \\
                              & Coupon frequency\footnote{$\ti := \{0.5, 1, \ldots, 4.5\}$} & Semi-annually \\
                              & Maturity            & $\T := 5$ \\
                              & Recovery            & $\gamma := 0\%$ \h
  \multirow{2}{*}{Put}        & Strike              & $\Kp := 105$ \\
                              & Period(s)           & $\Tp := \{3\}$ \h
  \multirow{2}{*}{Call}       & Strike\footnote{The strike of 110 is the clean price and needs to be adjusted for accrued interest, on a simple interest basis} & $\Kc := 110 + \C\tfrac{t \pmod {0.5}}{0.5}$ \\
                              & Period(s)           & $\Tc := [2, 5)$ \h
  \multirow{2}{*}{Conversion} & Quantity of stocks  & $\Kv := 1$ \\
                              & Period(s)           & $\Tv := [0, 5]$ \h
 \end{tabular}
 \caption{Convertible Bond Parameters}
 \label{tab:bond_parameters}
\end{table}
\end{savenotes}

\begin{savenotes}
\begin{table}[ht]
 \centering
 \begin{tabular}{|l|l|} \hline
  \theader{Parameter} & \theader{Value} \h
  Risk free rate      & $r := 5\%$ \h
  Volatility          & $\sigma := 20\%$ \h
  Hazard rate         & $\lambda := 2\%$ \h
  Default (total)     & $\eta := 100\%$ \h
  Default (typical\footnote{\citet{BP95} found that stock prices typically drop 30\% on announcement of default}) & $\eta := 30\%$ \h
  Default (partial)   & $\eta := 0\%$ \h
 \end{tabular}
 \caption{Stock Parameters}
 \label{tab:stock_parameters}
\end{table}
\end{savenotes}

% -----------------------------------------------------------------------------

\section{Sensitivity}
Figure \ref{fig:choices} illustrates the payoff of the convertible bond over it's lifetime, and the actions that led to the payoff (see Table \ref{tab:choices} for the list of actions, the associated payoff and the conditions required for the payoff).

\begin{figure}[htb!]
 \centering
 \includegraphics[width=14cm]{../common/fig_choices}
 \caption[Payoff]{Payoff of the convertible bond, with total default, and colours indicating the action taken for that payoff.  $\Dt = 2^{-3}$, $\DS = 2^{1}$ and $\S \in [0, 250]$.}
 \label{fig:choices}
\end{figure}

The figure has clear jumps at each coupon payment, producing a saw effect.  The other influencing factors are shown below:
\begin{itemize}
 \item \emph{Put:} the ``Put'' action is taken at time $\t = 3$ at stock prices from $\S = 0$ till $\S \approxeq 75$.  This option causes the value of the convertible bond prior to time $\t = 3$ to be elevated due to the higher payoff than the intrinsic value.
 \item \emph{Call:} the ``Call'' action is taken from time $\t = 2$ till almost $\t = 5$ with a narrow plateau around stock prices around $\S = 110$.  The call option does reduce the intrinsic value as apposed to being just the base for the forced conversion.
 \item \emph{Conversion:} on termination the ``Conversion'' action is taken for $\S \ge 104$.  At no other time is there voluntary conversion as, at time $\t = 2$ the upper payoff is similar to that of a forward and for similar reasons the value of the convertible bond increases.  Due to the increasing value the intrinsic value is never below the conversion threshold.
 \item \emph{Forced conversion:} the ``Forced conversion'' action is taken from time $\t = 2$ until almost $\t = 5$ above stock price $\S \approxeq 110$.
 \item \emph{Redemption:} on termination the ``Redemption'' action is taken for $\S < 104$.
 \item \emph{Hold:} the convertible bond is held for all other areas.  Of particular note is at time $\t = 3$ there is a narrow ``window'' from stock price $\S \approxeq 75$ to $\S \approxeq 110$ where the bond holder will not redeem, or be forced to redeem, the bond.  It is apparent that the changes of holding the bond beyond time $\t = 3$ is quite small.  Due to this the convertible bond value is dominated by the put, call and conversion options with the redemption value, and coupons after time $\t = 3$ having minimal impact.
\end{itemize}

\begin{figure}[htb!]
 \centering
 \includegraphics[width=14cm]{../common/fig_simple}
 \caption[Simple Call]{Comparison of initial value of the standard convertible bond, with total default, compared to one with a singular time for the call option.  Also the payoff of the ``simple call'' and colours indicating the action taken for the payoff.}
 \label{fig:simple}
\end{figure}\FloatBarrier

The above analysis suggests that the majority of impact that the call option has on the convertible bond's value is the leading edge at time $\t = 2$ and Figure \ref{fig:simple} illustrates the minimal change in price due to having a call with a single exercise time of $\Tc = \{2\}$.

In considering the sensitivity of the convertible bond to changes in each of the four components consider the following figures.

\begin{figure}[htb!]
 \centering
 \includegraphics[width=14cm]{../common/fig_varR}
 \caption[Varying Redemption Value]{Comparison of initial value of standard convertible bonds, with total default, to one with no redemption value.  Also the payoff of the ``no redemption'' and colours indicating the action taken for the payoff.}
 \label{fig:varR}
\end{figure}

\emph{Annuity}: Figure \ref{fig:varR} compares the convertible bond, with total default, with normal and no redemption value.  The payoff surface shows that at time $\t = 3$ the put effectively ``resets'' the payoff profile to that of a normal redemption value.

\begin{figure}[htb!]
 \centering
 \includegraphics[width=12cm]{../common/fig_varTp}
 \caption[Varying Put Time]{Comparison of initial value of convertible bonds, with total default, with different times for the put option.}
 \label{fig:varTp}
\end{figure}

\emph{Put:} Figure \ref{fig:varTp} compares the price of a convertible bond with differing put times.  The figure illustrates that the price of the bond is correlated with the put exercise time.  This effect happens because the coupon rate (converted to NACC) is greater than the risk free rate.  If the coupon rate were lower that an inverse correlation would be observed.  The effect dominant for $\S \in [0, 50]$ and tapers of until almost no impact by $\S = 100$.

\begin{figure}[htb!]
 \centering
 \includegraphics[width=12cm]{../common/fig_varTc}
 \caption[Varying Call Time]{Comparison of initial value of convertible bonds, with total default, with different times for the call option.}
 \label{fig:varTc}
\end{figure}

\emph{Call:} Figure \ref{fig:varTc} compares the price of a convertible bond with differing call times.  The figure illustrates that the price of the bond is correlated with the call exercise time.  This effect happens because the payoff behaves similarly to a forward on the stock and thus is increasing with time.

\begin{figure}[htb!]
 \centering
 \includegraphics[width=14cm]{../common/fig_varTv}
 \caption[Varying Call Time]{Comparison of initial value of convertible bonds, with total default, with different times for the conversion option.  Also the payoff of for $\Tv = [4, 5]$ and colours indicating the action taken for the payoff.}
 \label{fig:varTv}
\end{figure}

\emph{Conversion:} Figure \ref{fig:varTv} compares the price of a convertible bond with differing conversion times.  The figure illustrates the impact the conversion option, in conjunction with the call option, has on the price.  For conversion intervals that overlap with the leading edge of the call option (i.e. $\Tv = [1, 5]$ and $\Tv = [2, 4]$) the value is unchanged however when there is no overlap (i.e. $\Tv = [3, 5]$ and $\Tv = [4, 4]$) the upper section of the convertible bond's value is significantly suppressed limiting around the strike price of the call option.

As seen above the convertible bond is insensitive to changes in the redemption value and the majority of the value is determined by the leading edge of the call and put options (i.e. $\min(\Tc)$ and $\min(\Tp)$).  The call option influences the upper section, and the put option influences the lower section, of the pricing curve of the convertible bond.  Another important criteria is whether the conversion option on or before the leading edge of the call option or after the call option.

% -----------------------------------------------------------------------------

\section{Default Cost after Hedging}

\begin{figure}[htb!]
 \centering
 \includegraphics[width=14cm]{../common/fig_binTrade}
 \caption[Default Cost after Hedging]{Default Cost after Hedging, \H, for a convertible bond with typical default.}
 \label{fig:binTrade}
\end{figure}

Figure \ref{fig:binTrade} shows the un-adjusted default cost after hedging.  The shape shows areas where hedging was not possible: the upper right where the call and conversion (forced) options applied and the lower right where the put option limits the stock's ability to extend all the way down.  The hedging strategy is described in Equations \ref{eq:PDt} and \ref{eq:PDt2}.

The figure shows that the convertible bond has negative default cost after hedging.  If one were to synthesise the convertible bond then default events would be profitable due to a surplus (indicated by the negative in the graph) after paying the recovery component \X.

If one were to hedge a long position in the convertible bond then one would loose money in default due to the hedging strategy\footnote{In addition to the original loss due to default}.  However, the hedging will be profitable until default due to the ``insurance premium'' component of the hedging strategy.

The hedging strategy is not perfect however if a large enough ``population'' of uncorrelated convertible bonds are used in the hedging strategy then protection from hedging losses in default is possible.

% =============================================================================

\chapter{Conclusion}
This project derived and implemented two pricing methods of convertible bonds, that allowed for differing levels of default.  Parameters for three schemes, the explicit, implicit and Crank-Nicolson, were derived.

The binomial model was extended by considering the hedging strategy and showed that the binomial model does not provide for perfect hedging, requiring a cross financing of default and non-default events.  The cross financing requirement was shown to be \H, the unadjusted default cost after hedging, and a numerical analysis showed the convertible bond has a negative \H.

The sensitivities of the convertible bond was considered, with the put and call options influencing the lower and upper sections of the price curve respectively.  The conversion option can have impact on the price if it doesn't coincide with the leading edge of the call option.  The redemption value was shown to have little impact on the price.

Further avenues of research include:
\begin{itemize}
 \item considering the impact the bond and stock recovery has on the convertible bond price,
 \item exploration of better hedging strategies,
 \item Monte Carlo pricing, and
 \item simulations of ``insuring'' the convertible bond with respect to \H.
\end{itemize}

% =============================================================================

\appendix
\chapter{Appendix}
% -----------------------------------------------------------------------------

\section{Software}
All software using in this research project is open source.
\begin{savenotes}
\begin{table}[htb!]
 \centering
 \begin{tabular}{|l|l|} \hline
  \theader{Software} & \theader{Version} \h
  TeXLive    & 2012 \h
  Python     & 2.7.3 \h
  NumPy      & 1.6.2 \h
  SciPy      & 0.10.1 \h
  Matplotlib & 1.2.0 \h
  Git        & 1.7.11.5 \h
  pmake\footnote{As distributed with FreeBSD 9} & \h
 \end{tabular}
 \caption{Software}
\end{table}
\end{savenotes}

% -----------------------------------------------------------------------------

\section{Hardware}
\begin{table}[htb!]
 \centering
 \begin{tabular}{|l|l|} \hline
  \theader{Component} & \theader{Specification} \h
  CPU    & Intel Core i7 2600 \h
  Memory & 16GiB \h
 \end{tabular}
 \caption{Hardware}
\end{table}


% BIBLIOGRAPHY ---------------------------------------------------

\clearpage
\setlinespacing{1}
\addcontentsline{toc}{chapter}{Bibliography}
\bibliographystyle{../references/amsplnat}
\bibliography{../references/references.bib}
\end{document}
